\section{Introduction}

Latency-sensitive DNN inference increasingly relies on small fog clusters whose nodes differ in compute speed, memory, and communication quality. The scheduler must choose one node per request while respecting deadlines and resource limits. This discrete problem becomes difficult near congestion, where moving one task can change queueing and context-switching costs for many others \cite{Mao2017}.

Practical schedulers commonly separate candidate generation from feasibility enforcement. An optimizer first proposes a complete assignment, after which a repair operator reassigns or rejects tasks that violate memory, load, or deadline checks. This architecture is operationally useful, but it creates an objective mismatch: search is often ranked on raw assignments even though only repaired assignments are deployed. Our audit of the original pipeline showed that this mismatch can dominate the comparison. Raw GA-Hybrid and M-CA-ALM assignments agreed on only 23.1\% of task-node choices, whereas agreement rose to 60.8\% after shared repair. Consequently, a method may appear different internally while producing nearly identical service outcomes.

This paper addresses that mismatch directly. We propose an outcome-aligned version of the Memetic Congestion-Aware Augmented Lagrangian Method (M-CA-ALM). Candidate genotypes remain unrepaired to preserve population diversity, but elite selection evaluates the phenotype obtained after the same deterministic feasibility operator used at deployment. Augmented-Lagrangian local refinement is guarded: a refined candidate replaces its parent only when it improves a feasibility-first lexicographic criterion. The resulting search therefore optimizes the measured service outcome without collapsing the whole population into repaired copies.

The paper makes three contributions:
\begin{itemize}
    \item \textbf{Outcome-aligned memetic selection.} We introduce repair-aware phenotype ranking with non-Lamarckian genotype preservation, aligning candidate selection with final satisfied-task count while retaining structurally distinct candidates.
    \item \textbf{Guarded congestion-aware refinement.} We combine integer evolutionary exploration with sigmoid-relaxed ALM updates and accept a local update only when it improves feasibility first and satisfaction second.
    \item \textbf{A reproducible static-to-replay evaluation.} We fix the metric implementation and solver configuration, combine 30 paired static seeds with state-preserving replay, test repair-aware transfer in an independent GA family, and include a split-isolated offline graph policy. Paired tests, Holm correction, confidence intervals, and paired effect sizes are reported throughout.
\end{itemize}

The evidence supports a deliberately bounded claim. We call a load ``recoverable congestion'' when state-preserving replay retains a nonzero interval before the backlog reaches its cap. The ``immediate-overload boundary'' is the regime in which all tested controllers reach the cap at the same early slot. M-CA-ALM does not improve an already solved light-load regime. Its benefit appears in recoverable congestion, where it improves repaired satisfaction without sacrificing resource feasibility, at the cost of higher runtime. Formal replay confirms stronger service persistence at $N=180$. At the immediate-overload boundary at $N=300$, M-CA-ALM shows no reproducible advantage. Static experiments, independent transfer tests, and state-preserving replay are separate evidence tiers and are not pooled into one summary estimate. All conclusions are therefore limited to the fixed five-node simulator and tested load regimes; no cross-infrastructure, online-stability, capacity-planning, or hardware-deployment claim is made.
